\section{Ethical Considerations}
%
If deployed, our research will provide a tool for service providers to
safeguard against misuse of their services to store harmful
content. By introducing independent auditors, we reduce the potential
risk of surveillance and censorship. However, the ethical risks of
general content moderation itself needs to be taken into account when
evaluating the ethical implications of our work. The debate regarding
free speech vs.\ moderating misinformation does not have a conclusive
answer but several works have tried to summarize arguments both in
favor or against it~\cite{kozyreva2023resolving}.

Nonetheless, service providers have a responsibility to ensure that
their services are not being misused to proliferate harmful
content. Therefore, moderating for such content is justified, and
arguably necessary~\cite{howard2023ethics}. However, to protect the
interests of the clients, we need to ensure that auditors of the
blocklist act independently and without influence. While in many cases
this may not be possible, e.g., due to government interference, we
believe this trust model is more suitable than endowing trust entirely
on self-interested service providers.

All that said, our research itself, as a design for a content-moderation system, poses little ethical risk, at least directly.  (The parameters and policies around its deployment need to be carefully considered, however.)  Note, moreover, that our research
did not involve testing with any real-world data or human subjects. As
such, no direct harms or violations of privacy resulted from this
research. We acknowledge that, sometimes, the most ethical path is not to do the
research or not to publish the research after it is complete. However,
due to the urgent need to combat the proliferation of CSAM, we believe
that the potential benefits of our research far outweigh its risks,
thereby satisfying the Beneficence criterion of the Menlo Report.
